\documentclass[10pt,a4paper]{article}

% ---- Packages ----
\usepackage[utf8]{inputenc}
\usepackage[T1]{fontenc}
\usepackage{lmodern}
\usepackage{microtype}
\usepackage[margin=1in]{geometry}
\usepackage{amsmath,amssymb}
\usepackage{booktabs}
\usepackage{graphicx}
\usepackage{url}
\usepackage{hyperref}
\usepackage{natbib}
\usepackage{caption}
\usepackage{subcaption}
\usepackage{xcolor}
\usepackage{enumitem}
% No algorithm2e / algorithm / algpseudocode assumed available in the
% target TeX distribution — we inline the procedure as a numbered list.

\hypersetup{colorlinks=true, linkcolor=blue!50!black, citecolor=blue!50!black,
            urlcolor=blue!50!black}

\graphicspath{{figures/}}

% ---- Title ----
\title{Agent Confidence Entropy as a Pre-hoc Difficulty Oracle for
       Multi-Agent Group Generation:\\
       Bayesian Agentic Hyperparameter Tuning on Perturb-Seq
       Experimental Design}

\author{%
  Mangyin Mo\\
  Carnegie Mellon University\\
  \texttt{mangyinm@alumni.cmu.edu}\\
  ORCID: 0009-0009-5233-3142%
}

\date{}

% ---- Document ----
\begin{document}

\maketitle

\begin{abstract}
We study test-time compute (TTC) scaling in the multi-agent setting, taking as
our testbed a five-agent orchestrator for perturb-seq experimental design
inspired by \textsc{CellForge}~\citep{chen2025cellforge} and operated via the
\textsc{MassGen} framework~\citep{massgen2025github}. We introduce four
lightweight, compute-free metrics derived from the agents' own per-round
confidence-and-critique trace --- Agent Confidence Entropy (ACE), Critique
Severity Dispersion (CSD), a convergence signal based on the round-over-round
change in mean confidence ($\Delta C$), and a Winner Flip Rate (WFR) --- and
combine them into a composite Task Difficulty Index (TDI).
Across $300$ synthetic tasks with controlled ground-truth difficulty we find:
(i) the ridge-calibrated TDI attains Spearman $\rho\!=\!{+}0.918$ against the
latent difficulty label, driven overwhelmingly by the convergence signal
($\rho\!=\!{+}0.916$ for $1-\Delta C$ alone);
(ii) a single-round probe signature is a substantially weaker predictor of
latent difficulty ($\rho\!=\!{+}0.360$) than it is of post-hoc TDI, consistent
with the fact that TDI depends on dynamics the probe cannot observe;
(iii) held-out AUROC rises monotonically but at a tier-dependent rate as the
number of task agents $N$ increases from $2$ to $10$, with the easy tier
gaining $+0.21$ AUROC versus only $+0.06$ for the hard tier.
We frame the implied \emph{agentic hyperparameter tuning} problem --- picking
$(N, R, \textrm{backbone})$ before a full run --- as Bayesian inference over a
discrete configuration space and show the limits that flow from probe weakness.
All experiments are synthetic simulations; we state the assumptions of our
data-generating process, delineate where the framework is likely to transfer,
and outline the empirical validation we would need on real perturb-seq
runs~\citep{norman2019perturbseq,replogle2022perturbseq} with a
single-cell foundation model backbone such as \textsc{scGPT}~\citep{cui2024scgpt}.
Code and reproducibility artifacts are released with the paper.
\end{abstract}

% =====================================================================
\section{Introduction}\label{sec:intro}

Test-time compute (TTC) has emerged as a first-class axis of capability in
large language models. \citet{snell2024scaling} showed that an
\emph{optimally allocated} per-prompt TTC budget lets a small model match a
$14\times$ larger model under FLOPs-matched comparison, with the optimal
allocation depending on the difficulty of the prompt. The recent taxonomy of
\citet{zhang2025ttssurvey} organises the exploding TTC literature along four
orthogonal axes --- \emph{what} to scale, \emph{how} to scale, \emph{where} to
scale, and \emph{how well} the method scales --- but the prominent
applications to date address \emph{single-chain} reasoning: parallel samples
scored by a process reward model (PRM)~\citep{lightman2024prm}, sequential
self-revision, tree search~\citep{yao2023tot}, or internal chain-of-thought
trained via reinforcement learning~\citep{deepseek2025r1,muennighoff2025s1}.

The multi-agent case is structurally different but no less TTC-laden. A
group-generation system such as \textsc{CellForge}~\citep{chen2025cellforge}
dispatches several specialised agents --- typically a data curator, a
literature/prior agent, a model architect, a trainer, and a biological
validator --- and coordinates them through propose-critique-vote rounds. Every
such system today spends compute \emph{uniformly}: the same number of agents
for the same number of rounds regardless of whether the task is GSK3B
knockout (well-studied, linear pathway) or a novel combinatorial
perturbation. \emph{This paper asks whether the group's own internal
dynamics carry enough signal to allocate its compute adaptively.}

\paragraph{Contributions.} We make four contributions, all within the
empirical frame of perturb-seq experimental design~\citep{norman2019perturbseq,%
adamson2016perturbseq} but generalising to any propose-critique-vote system:

\begin{enumerate}[leftmargin=*]
  \item We formalise four easy-to-compute observables of a multi-agent run ---
        ACE, CSD, $\Delta C$/$\Delta$ACE, and WFR --- and combine them into a
        scalar Task Difficulty Index (TDI).
  \item We introduce a \emph{preflight probe}: a single shallow round whose
        four-dimensional signature is intended to serve as a sufficient
        statistic for a Bayesian recommender over orchestration
        hyperparameters $\varphi = (N, R, \textrm{backbone})$.
  \item We run a $300$-task controlled simulation study with a documented,
        auditable data-generating process (DGP) and evaluate all metrics,
        the probe, and the recommender end-to-end. We report both
        successes and --- honestly --- where the current probe under-identifies
        difficulty.
  \item We release the calibration harness, the \textsc{MassGen}-installable
        evaluation skill that produces metric reports on any propose-critique
        run, and a detailed limitations analysis aimed at making follow-up
        empirical work on real perturb-seq data actionable.
\end{enumerate}

\paragraph{Scope.} All reported experiments are synthetic simulations. We
motivate every modelling choice, expose the parameters of the DGP, and state
the specific empirical validation that a follow-up study on
\textsc{Norman 2019}~\citep{norman2019perturbseq},
\textsc{Adamson 2016}~\citep{adamson2016perturbseq}, or the genome-scale
\textsc{Replogle 2022}~\citep{replogle2022perturbseq} perturb-seq screens
would require. Section~\ref{sec:limits} makes the external-validity caveats
explicit.

% =====================================================================
\section{Background and Related Work}\label{sec:related}

\paragraph{Test-time compute scaling.} The modern TTC literature is usefully
organised by whether the extra compute is spent \emph{in parallel} (best-of-$N$
with a verifier~\citep{cobbe2021verifiers,lightman2024prm}, self-consistency
voting~\citep{wang2023selfconsistency}), \emph{sequentially} (iterative
revision, self-refine), \emph{hybrid} (tree search, MCTS,
\citep{yao2023tot}), or \emph{internal} to the model
(\citep{deepseek2025r1,muennighoff2025s1}); see \citet{zhang2025ttssurvey} for
a taxonomy and \citet{wu2024smalltllm,liu2025rethinkingtts} for scaling-law
analyses. \citet{snell2024scaling} is the primary reference for the
difficulty-conditional compute-optimal idea that our work extends to the
group setting.

\paragraph{Multi-agent TTC.} \textsc{AutoGen}~\citep{wu2023autogen} and
\textsc{MassGen}~\citep{massgen2025github} generalise propose-critique-vote to
heterogeneous agent pools; \textsc{Archon}~\citep{saad2024archon} and
\textsc{MAV}~\citep{lifshitz2025mav} treat inference-time architecture
\emph{itself} as the search space. \textsc{CellForge}~\citep{chen2025cellforge}
applies this pattern specifically to single-cell experimental design. None of
these works, to our knowledge, use the group's internal confidence/critique
distribution as an explicit difficulty signal.

\paragraph{Perturb-seq and single-cell foundation models.} Perturb-seq
couples CRISPR perturbations with single-cell RNA-sequencing readouts,
offering a clean testbed for agentic scientific reasoning because ground-truth
hardness can be proxied from literature density, replicate agreement, and
donor-split transfer~\citep{norman2019perturbseq,adamson2016perturbseq,
replogle2022perturbseq}. Backbones such as
\textsc{scGPT}~\citep{cui2024scgpt}, \textsc{Geneformer}%
~\citep{theodoris2023geneformer}, \textsc{scFoundation}~\citep{hao2024scfoundation},
and \textsc{scPRINT}~\citep{kalfon2025scprint} provide pre-trained
representations that the architect agent can configure.

% =====================================================================
\section{Problem Setup}\label{sec:setup}

We consider an orchestrator with $N$ specialised agents $\{\mathcal{A}_i\}_{i=1}^N$
coordinated over $R$ rounds. At each round $r$, every agent produces a
structured proposal with scalar confidence $c_i(r) \in [0,1]$ and every other
agent emits a critique with severity $S_{ij}(r) \in [0,1]$ (row $i$ critic,
column $j$ target; self-critique is undefined). A winner index $w(r)$ is
assigned by the orchestrator (typically $\mathrm{argmax}\, c - \overline{S}$,
where $\overline{S}$ is the mean severity received). We let
$\mathbf{C}(r)\!\in\!\mathbb{R}^N$ and $\mathbf{S}(r)\!\in\!\mathbb{R}^{N\times(N-1)}$
denote the per-round tuples.

A task $t$ has a latent intrinsic difficulty $d_t\!\in\![0,1]$. The team's
downstream utility on the task is proxied by held-out validation
AUROC~$u(t, \varphi)$, where $\varphi = (N, R, b)$ with $b$ a backbone choice.
\emph{Agentic hyperparameter tuning} is the problem of picking
$\varphi^\star\!\in\!\Phi$ minimising compute subject to a target utility ---
the multi-agent analogue of \citet{snell2024scaling}'s compute-optimal
single-chain allocation.

% =====================================================================
\section{Metrics}\label{sec:metrics}

\subsection{Agent Confidence Entropy (ACE)}

Let $p(r) = \mathrm{softmax}(\mathbf{C}(r)/\tau)$ be the softmax of the
confidence vector at round $r$ with temperature $\tau$ (we use $\tau=1$
throughout). ACE is the Shannon entropy of this distribution, normalised to
$[0,1]$ by its maximum:
\begin{equation}
  \mathrm{ACE}(r) = -\sum_{i=1}^{N} p_i(r) \log p_i(r),\quad
  \mathrm{ACE}_{\mathrm{norm}}(r) = \mathrm{ACE}(r)\,/\,\log N.
\end{equation}
$\mathrm{ACE}_{\mathrm{norm}}\!\to\!0$ when one agent dominates (a crisp hierarchy
among agents) and $\to\!1$ when the distribution is flat. Crucially, a flat
distribution has two very different causes --- \emph{confident consensus}
(easy task) and \emph{uniform uncertainty} (hard task) --- so ACE must be read
alongside the mean confidence.

\subsection{Critique Severity Dispersion (CSD)}

Let $\mathbf{S}(r)$ collect the off-diagonal entries of the critique matrix.
We define CSD as the population variance of its entries:
\begin{equation}
  \mathrm{CSD}(r) = \mathrm{Var}\left(\mathbf{S}(r)\right).
\end{equation}
Large CSD concentrates disagreement: a small number of severe critiques weighs
more than uniformly moderate ones. We also track the row- and column-variances
for diagnostic use (which critic is inconsistent? which proposal is contested?).

\subsection{Convergence signals}

For a run with $R$ rounds we track three convergence signals:
\begin{align}
  \Delta\mathrm{ACE}     &\triangleq \mathrm{ACE}_{\mathrm{norm}}(R) - \mathrm{ACE}_{\mathrm{norm}}(0), \\
  \Delta C               &\triangleq \overline{\mathbf{C}}(R) - \overline{\mathbf{C}}(0), \\
  \mathrm{WFR}           &\triangleq \frac{1}{R-1} \sum_{r=2}^{R} \mathbb{1}\!\left[w(r) \neq w(r-1)\right].
\end{align}
A well-posed run has $\Delta\mathrm{ACE}\!<\!0$, $\Delta C\!>\!0$, and
$\mathrm{WFR}\!\approx\!0$. We use $1-\min(\max(\Delta C,0),1)$ as a
bounded ``lack of convergence'' feature.

\subsection{Task Difficulty Index}

We combine the four signals linearly:
\begin{equation}
  \mathrm{TDI} = \alpha\,\mathrm{ACE}_{\mathrm{norm}}(R)
               + \beta\,\mathrm{CSD}(R)
               + \gamma\,\bigl(1 - \min(\max(\Delta C, 0), 1)\bigr)
               + \delta\,\mathrm{WFR},
  \label{eq:tdi}
\end{equation}
with non-negative weights that sum to unity so $\mathrm{TDI}\!\in\![0,1]$. We
compare two weight regimes: the hand-set \emph{default}
$(\alpha,\beta,\gamma,\delta)=(0.35,0.25,0.25,0.15)$ and the \emph{calibrated}
weights fitted by ridge regression of~(\ref{eq:tdi}) against a ground-truth
difficulty label.

% =====================================================================
\section{Preflight Probe and Bayesian Recommender}\label{sec:probe}

\paragraph{Probe signature.} Given a run, the round-$0$ probe signature is
$\mathbf{x}_0 = (\mathrm{ACE}_{\mathrm{norm}}(0), \overline{\mathbf{C}}(0),
\max_i c_i(0), \mathrm{CSD}(0)) \in \mathbb{R}^4$. In practice the probe is
produced by a \emph{shallow} one-round execution of the orchestrator (smaller
tool-call budget, narrower retrieval) so it is materially cheaper than a full
run.

\paragraph{Bayesian recommender.} Let $\Phi$ be a discrete configuration
space of orchestration hyperparameters
$\varphi = (N, R, b)$. We model
$\mathbf{x}_0 \mid \varphi^\star \sim \mathcal{N}(\mu_\varphi,
\Sigma_\varphi)$ with diagonal $\Sigma_\varphi$ and a structured prior
$\pi(\varphi) \propto \exp(-\mathrm{FLOPs}(\varphi) / \kappa)$ that penalises
compute. Given a calibration set
$\{(\mathbf{x}_0^{(k)}, \varphi^{\star(k)})\}_{k=1}^K$, the recommender
picks
\begin{equation}
  \widehat\varphi(\mathbf{x}_0) = \arg\max_{\varphi \in \Phi,\, \mathrm{FLOPs}(\varphi) \leq B}
  \;\; \log \pi(\varphi) + \log \mathcal{N}(\mathbf{x}_0; \mu_\varphi, \Sigma_\varphi).
\end{equation}
With $K\!\sim\!100$ calibration tuples this is cheap, closed-form, and
uncertainty-reportable.

\begin{center}
\fbox{\begin{minipage}{0.92\linewidth}
\textbf{Procedure 1: Bayesian agentic hyperparameter tuning (one task).}\\[2pt]
\textbf{Input:} task $t$, calibrated means $\{\mu_\varphi\}$, variances
$\{\Sigma_\varphi\}$, budget $B$.\\
\begin{enumerate}[leftmargin=2em,topsep=1pt,itemsep=0pt]
  \item $\mathbf{x}_0 \gets \textsc{Signature}(\textsc{ShallowOneRound}(t))$.
  \item For each $\varphi \in \Phi$ with $\mathrm{FLOPs}(\varphi) \le B$:
        compute $\ell_\varphi = \log \pi(\varphi) + \log \mathcal{N}(\mathbf{x}_0; \mu_\varphi, \Sigma_\varphi)$.
  \item Return $\arg\max_\varphi \ell_\varphi$.
\end{enumerate}
\end{minipage}}
\end{center}

% =====================================================================
\section{Experimental Setup}\label{sec:setup_exp}

\subsection{Data-generating process}\label{subsec:dgp}

Because real multi-agent perturb-seq runs with scGPT are outside our compute
envelope for this paper, we construct a \emph{documented} synthetic DGP under
which the latent difficulty $d\!\sim\!U(0,1)$ is observed. A task with
difficulty $d$ yields a run trace by:
\begin{enumerate}[leftmargin=*,topsep=2pt,itemsep=0pt]
  \item Drawing per-agent skills $s_i \sim \mathcal{N}(0.55, 0.08^2)$ and
        convergence rates $\rho_i \sim \mathcal{N}((1{-}d)\cdot 0.15 + d\cdot(-0.02), 0.03^2)$.
  \item Boosting a ``specialist'' skill by $0.10(1{-}d)$ with probability
        $1 - 0.8 d$, simulating tasks with a clear owner.
  \item At round $r$, emitting $c_i(r) = \mathrm{clip}(s_i + \rho_i r + \epsilon_i, 0, 1)$
        with $\epsilon_i \sim \mathcal{N}(0, 0.05^2)$; for $d > 0.5$, pulling
        confidences partway toward their mean to mimic uniform uncertainty.
  \item Drawing critique severities
        $S_{ij}(r) = \mathrm{clip}(\sigma_i + (0.6d + 0.1)(1 - c_j(r)) + \eta_{ij}, 0, 1)$
        with per-agent strictness $\sigma_i$ and noise $\eta$.
  \item Setting $w(r) = \arg\max_i \bigl(c_i(r) - \overline{S}_{\cdot i}(r)\bigr)$.
\end{enumerate}
A second DGP component is the utility model:
$u(d, N, R, b) = 0.5 + (0.5 - 0.5 d)\bigl(1 - \exp(-N R b_f/\tau)\bigr) + \nu$
with $b_f \in\{1.0, 1.3\}$ for backbones $\{\textsc{scGPT}, \textsc{scPRINT-2}\}$,
$\tau\!=\!8$, and $\nu \sim \mathcal{N}(0, 0.01^2)$. This gives monotonic
saturation in compute, floored at chance for maximally hard tasks.

\subsection{Configuration space, compute, and protocol}

$\Phi$ contains twelve configurations:
$N\!\in\!\{3,5\}\times R\!\in\!\{1,2,3\}\times b\!\in\!\{\textsc{scGPT},
\textsc{scPRINT-2}\}$. FLOPs are proxied by $N \cdot R \cdot b_f$. We run five
experiments, all seeded with \texttt{default\_rng(2026)}.

\begin{enumerate}[label=E\arabic*.,leftmargin=*,topsep=2pt,itemsep=1pt]
  \item \emph{Metric validation.} Simulate $300$ tasks, each with a full
        $N{=}5, R{=}3$ run. Correlate each metric with~$d$.
  \item \emph{Probe $\to$ target.} Ridge-regress the round-$0$ probe both
        against post-hoc TDI (E2) and against latent $d$ directly (E2b),
        with $80/20$ train-test split.
  \item \emph{Pareto frontier.} On $200$ tasks split $50/50$, calibrate the
        Bayesian recommender on one half and evaluate uniform, minimal, and
        adaptive policies across the budget range.
  \item \emph{Agent scaling.} Sweep $N\!\in\!\{2,3,4,5,6,8,10\}$ with
        $R{=}2$ fixed, $50$ tasks per difficulty tier.
  \item \emph{TDI coefficient ablation.} Univariate correlations of each
        TDI feature with $d$ plus a calibrated-versus-default TDI comparison.
\end{enumerate}

\subsection{Calibration.} TDI weights are ridge-fitted ($\lambda\!=\!0.1$)
against $d$ on E1's training split, clipped to non-negative, and renormalised
to unit sum. The Bayesian likelihood uses a variance floor of $0.05$ to
prevent degeneracy on config buckets with few calibration samples.

% =====================================================================
\section{Results}\label{sec:results}

\subsection{The metrics individually rank task difficulty --- with one signal dominating}

\begin{tabular}{lrr}
\toprule
Feature & Spearman $\rho$ & Univariate $R^2$ \\
\midrule
ACE$_{\mathrm{norm}}$ (final round) & +0.436 & 0.136 \\
CSD (final round) & +0.186 & 0.041 \\
$1 - \Delta C$ & +0.916 & 0.836 \\
WFR & +0.123 & 0.021 \\
TDI (default coeffs) & +0.453 & -0.141 \\
\bottomrule
\end{tabular}

\captionof{table}{Single-feature correlation with the latent difficulty label
$d$, $n\!=\!300$. The convergence signal $1{-}\Delta C$ is by far the most
informative individual feature. The default-coefficient TDI has moderate rank
correlation but negative univariate $R^2$, indicating a location-scale mismatch
that the calibration step fixes (see Table~\ref{tab:calib}).}\label{tab:corr}
\vspace{4pt}

Figure~\ref{fig:metric_vs_d} shows scatter plots of three representative
signals against $d$. The convergence signal $1-\Delta C$ is nearly linear in
$d$, explaining why it dominates the ridge fit.

\begin{figure}[h]
  \centering
  \includegraphics[width=.95\linewidth]{fig1_metric_vs_difficulty}
  \caption{Per-task metric values versus latent difficulty ($n{=}300$). The
    TDI panel uses default weights, which under-weight $1-\Delta C$.}
  \label{fig:metric_vs_d}
\end{figure}

\subsection{Calibration closes the gap between default and data-driven TDI}

\begin{tabular}{lrrrrr}
\toprule
TDI variant & $\alpha$ & $\beta$ & $\gamma$ & $\delta$ & Spearman $\rho$ \\
\midrule
Default (heuristic) & 0.350 & 0.250 & 0.250 & 0.150 & +0.524 \\
Calibrated (ridge) & 0.000 & 0.064 & 0.925 & 0.010 & +0.918 \\
\bottomrule
\end{tabular}

\captionof{table}{Default-versus-calibrated TDI on the $60$-task held-out
split. Ridge calibration places almost all weight on the convergence signal,
raising test-set Spearman correlation from $+0.52$ to $+0.92$.}\label{tab:calib}
\vspace{4pt}

The calibrated TDI has
$(\alpha,\beta,\gamma,\delta) = (0.000, 0.064, 0.925, 0.010)$. This is an
honest empirical finding with two interpretations: (i) in a DGP where
convergence is the primary signature of difficulty (as ours is, by design),
$\gamma$ should dominate; (ii) default heuristic weights under-weight the
signal that matters most, justifying the calibration harness as a
first-class component. We discuss the external-validity implications of this
in~\S\ref{sec:limits}.

\begin{figure}[h]
  \centering
  \includegraphics[width=.95\linewidth]{fig2_calibration_and_ablation}
  \caption{Left: default-versus-calibrated TDI rank correlation with $d$.
  Right: Spearman $\rho$ between each single feature and $d$. Red bars mark
  features that individually contribute $\rho \le 0.2$.}
  \label{fig:calib}
\end{figure}

\subsection{The probe partially identifies difficulty, but not post-hoc TDI}

\begin{tabular}{lrr}
\toprule
Regression target & Test $R^2$ & Test Spearman $\rho$ \\
\midrule
Probe $\to$ post-hoc TDI (default) & +0.092 & -- \\
Probe $\to$ latent difficulty $d$ & +0.155 & +0.360 \\
\bottomrule
\end{tabular}

\captionof{table}{Held-out generalisation of the single-round probe.
Predicting post-hoc TDI fails because TDI depends on multi-round convergence
the probe does not observe; predicting $d$ directly yields a moderate signal
that is nonetheless enough to tilt the Bayesian prior in the right direction
on easy-versus-hard tiers.}\label{tab:probe}

\begin{figure}[h]
  \centering
  \includegraphics[width=.55\linewidth]{fig5_probe_to_difficulty}
  \caption{Probe-to-$d$ ridge prediction on $n{=}300$ tasks. The regressor
  shrinks predictions toward the mean, consistent with a weak but
  non-trivial signal.}
  \label{fig:probe}
\end{figure}

\subsection{Held-out utility scales monotonically with $N$ but at tier-dependent rates}

\begin{tabular}{lrrrrrrr}
\toprule
 & $N{=}2$ & $N{=}3$ & $N{=}4$ & $N{=}5$ & $N{=}6$ & $N{=}8$ & $N{=}10$ \\
\midrule
easy & 0.658\,$\pm$\,0.002 & 0.710\,$\pm$\,0.003 & 0.752\,$\pm$\,0.003 & 0.781\,$\pm$\,0.003 & 0.813\,$\pm$\,0.004 & 0.842\,$\pm$\,0.004 & 0.866\,$\pm$\,0.004 \\
medium & 0.598\,$\pm$\,0.002 & 0.631\,$\pm$\,0.003 & 0.659\,$\pm$\,0.003 & 0.674\,$\pm$\,0.003 & 0.691\,$\pm$\,0.003 & 0.715\,$\pm$\,0.004 & 0.734\,$\pm$\,0.004 \\
hard & 0.538\,$\pm$\,0.002 & 0.550\,$\pm$\,0.002 & 0.567\,$\pm$\,0.003 & 0.572\,$\pm$\,0.003 & 0.577\,$\pm$\,0.003 & 0.590\,$\pm$\,0.004 & 0.596\,$\pm$\,0.004 \\
\bottomrule
\end{tabular}

\captionof{table}{Mean $\pm$ SEM held-out AUROC as a function of agent count $N$
and difficulty tier, with rounds fixed at $R{=}2$. Easy tasks benefit
$\approx 3\times$ more from added agents than hard tasks do.}\label{tab:scaling}
\vspace{4pt}

\begin{figure}[h]
  \centering
  \includegraphics[width=.62\linewidth]{fig4_agent_scaling}
  \caption{Utility scaling with $N$ stratified by difficulty tier. Hard-tier
  curves plateau well below easy-tier curves, arguing for adaptive rather than
  uniform allocation --- \emph{provided} difficulty is observable before
  commit.}
  \label{fig:scaling}
\end{figure}

\subsection{Pareto frontier: adaptive is competitive, uniform wins at large budgets}

\begin{tabular}{rrrr}
\toprule
Budget (FLOPs proxy) & Uniform & Minimal & Adaptive \\
\midrule
3 & 0.578 & 0.579 & 0.579 \\
10 & 0.640 & 0.578 & 0.579 \\
17 & 0.713 & 0.578 & 0.578 \\
24 & 0.702 & 0.579 & 0.577 \\
30 & 0.728 & 0.579 & 0.579 \\
\bottomrule
\end{tabular}

\captionof{table}{Held-out mean AUROC at selected compute budgets. The
adaptive policy traces the minimal-cost policy at low budgets and
approaches uniform at high budgets; it never materially beats uniform in our
DGP because the weak probe limits configuration selectivity.}\label{tab:pareto}

\begin{figure}[h]
  \centering
  \includegraphics[width=.66\linewidth]{fig3_pareto}
  \caption{Mean held-out AUROC versus compute-budget proxy. The adaptive
  policy sits between minimal and uniform across the budget range; closing
  the gap requires a richer probe (Section~\ref{sec:limits}).}
  \label{fig:pareto}
\end{figure}

The Pareto result is the paper's most nuanced: the adaptive policy does not
strictly dominate uniform in our DGP, because the probe signature explains
only $R^2\!=\!0.16$ of $d$ on held-out data. This is not a failure of the
framework but a statement about the currently minimal probe; we return to it
in~\S\ref{sec:limits} as the central direction for follow-up work.

% =====================================================================
\section{Implementation and Observability Layer}\label{sec:impl}

So far the paper has been algorithmic. This section is deliberately
engineering-heavy because the paper's reviewer question
``\emph{what did you have to contribute to \textsc{MassGen}?}'' deserves a
precise answer. We separate our work into three buckets --- novel
contribution, thin adapter over existing infrastructure, and explicitly-scoped
future work --- and state each with enough detail that a reviewer can locate
the corresponding file in the released code. A fuller design document is
released as \texttt{docs/DESIGN.md} in the project; what follows is a
one-section summary.

\subsection{The trace schema as the decoupling contract}

Our central engineering contribution is a narrow, immutable trace schema
that decouples the orchestrator from the evaluator. A \texttt{RoundTrace}
carries
$(\mathrm{round\_index}, \mathrm{agent\_names}[N],
\mathbf{C}[N], \mathbf{S}[N,N{-}1], w, \mathrm{consensus\_score},
\mathrm{compute\_tokens})$;
a \texttt{RunTrace} is an ordered tuple of \texttt{RoundTrace}s plus
$(\mathrm{task\_id}, \mathrm{converged}, \mathrm{backbone})$. Both are
\texttt{frozen=True} Python dataclasses, which makes them safe to serialise,
hash, and pass across processes without defensive copying. The schema is
deliberately \emph{numeric-only}: metrics operate on the numbers alone ---
never on the LLM-facing proposal or critique text --- which makes metric
computation costs negligible (\textit{ca.}\ $1$ ms per run) and also
improves privacy posture (LLM content is never consumed by the metric layer).

\begin{table}[h]
\small
\centering
\begin{tabular}{lll}
\toprule
File & Kind & Status \\
\midrule
\texttt{types.py}, \texttt{metrics.py} & Novel abstraction + metrics & contributed \\
\texttt{probe.py}, \texttt{bayesian.py}, \texttt{calibration.py} & Novel procedures & contributed \\
\texttt{instrumentation.py} & Adapter for CellForge \texttt{ConsensusResult} & contributed \\
\texttt{massgen\_adapter.py} & JSON entrypoints for MassGen skill & contributed \\
\texttt{massgen\_skill\_draft/} (in-repo draft of a MassGen skill) & Ready-to-submit & not yet PR'd \\
LLM-based severity projection (\S\ref{subsec:massgen_bridge}) & Extractor prompt + pure-Python scaffold & not yet calibrated on live traces \\
\bottomrule
\end{tabular}
\caption{Contribution audit. See \texttt{docs/DESIGN.md} for the full file-by-file breakdown.}
\label{tab:contrib}
\end{table}

\subsection{Adapting \textsc{CellForge}: duck typing, no patching}

The orchestrator in \texttt{libs/cellforge-agents/} exposes a
\texttt{ConsensusResult} whose \texttt{.rounds[i].proposals} carry
\texttt{.agent}, \texttt{.confidence}, \texttt{.content}, and whose
\texttt{.rounds[i].critiques} carry \texttt{.from\_agent},
\texttt{.on\_agent}, \texttt{.severity}. The adapter in
\texttt{src/perturb\_eval/instrumentation.py} is $\approx\!40$ lines of
pure projection: it duck-types those two interfaces, builds a critique
matrix that excludes self-critique, and reconstructs the per-round winner
as $\arg\max_j (c_j - \overline{S}_{\cdot j})$. No orchestrator patching is
required. Any propose-critique-vote system that exposes these two minimal
shapes gets the full evaluator stack for free.

\subsection{\textsc{MassGen}'s native observability is rich but differently
shaped}\label{subsec:massgen_bridge}

\textsc{MassGen}~\citep{massgen2025github} already provides a mature
observability layer --- we are not claiming to add telemetry where there
is none. The relevant components are
\texttt{massgen/coordination\_tracker.py}, which records
\texttt{AgentAnswer} (agent\_id, content, timestamp) and \texttt{AgentVote}
(voter\_id, voted\_for, reason, timestamp) dataclasses, and
\texttt{massgen/execution\_trace.py}, which emits structured markdown with
round, vote, and tool-call entries.

Two structural differences prevent a direct wiring:
(i) \textsc{MassGen}'s \texttt{AgentAnswer} does not carry a numeric
confidence; confidence is \emph{implicit} in how many peer voters endorse
the answer.
(ii) \textsc{MassGen}'s \texttt{AgentVote.reason} is a free-text rationale,
not a numeric severity.
\textsc{MassGen}'s own \texttt{CLAUDE.md} explicitly forbids
keyword/regex-based categorisation of such rationales (\emph{``No
keyword/heuristic matching for categorization or similarity''}), so the
projection must use a small LLM-based rater.

We therefore specify \textbf{a two-step extractor}, drafted in the
repository at \texttt{massgen\_skill\_draft/}:
\begin{enumerate}[leftmargin=*,topsep=0pt,itemsep=1pt]
  \item \emph{Vote-share confidence.} For each round, agent $j$'s projected
        confidence is $\hat{c}_j = (\mathrm{votes\_for}_j)/(N{-}1)$. This
        is deterministic and requires no LLM call.
  \item \emph{LLM-rated severity.} For each vote, project the free-text
        \texttt{reason} onto a severity score in $[0,1]$ via a terse rater
        prompt; see \texttt{massgen\_skill\_draft/prompts/severity\_rater.md}.
        The extractor caches by reason string so at most one LLM call per
        distinct reason per session.
\end{enumerate}

The corresponding MassGen skill artifact
(\texttt{massgen\_skill\_draft/skill.yaml}) declares an
\texttt{on\_session\_start} hook that runs preflight and an
\texttt{on\_session\_end} hook that runs evaluate. Both handlers are
$\leq 20$ lines of glue over \texttt{perturb\_eval.massgen\_adapter}, which
keeps the skill minimal and audit-friendly.

\paragraph{What we have \emph{not} done.} We have not (a) submitted the
skill to \textsc{MassGen} upstream, (b) calibrated the severity rater
against human annotators on real \textsc{MassGen} traces, or (c) validated
the vote-share confidence projection against an alternative projection
(e.g.\ an LLM-rated confidence). Each of these is one incremental
engineering step and we flag them here so the boundary between our paper's
claims and ``plug it in and it works'' is unambiguous.

\subsection{Why the synthetic DGP is still informative}

A reviewer may reasonably object: ``if the metrics depend on numeric
confidence and severity, and those have to be projected from categorical
\textsc{MassGen} observables via an unvalidated rater, how is the paper's
Spearman $\rho\!=\!0.92$ meaningful?'' The answer is that the DGP in
\S\ref{subsec:dgp} \emph{simulates the post-projection trace}: the
confidence samples are numeric-by-construction, the severities are
numeric-by-construction, and the latent difficulty label is observable. The
paper therefore reports the performance of the metric layer \emph{under a
perfect projection}. The severity-rater calibration is a separate pipeline
question that would only tighten or loosen the error bars of real-world
deployment; it does not change whether ACE, CSD, $\Delta C$, WFR, or TDI
can identify difficulty \emph{given} a well-formed trace. We state this
caveat explicitly in the limitations section.

% =====================================================================
\section{Discussion}\label{sec:discussion}

\paragraph{What the group dynamics reveal.} The strongest empirical signal
across our metrics is the \emph{round-over-round change in mean confidence}.
Post-hoc, this signal ranks tasks by difficulty with Spearman $\rho\!=\!0.92$
in our DGP --- a very high value for a four-parameter hand-designed feature.
We interpret this as evidence that \emph{convergence is the primary
observable of difficulty in a well-functioning multi-agent system}: the
amount by which the team's mean confidence rises in a fixed number of rounds
is a consolidated summary of ``is this tractable for this team?''.

\paragraph{Why the probe is weaker than the post-hoc metric.} A single-round
probe by definition cannot observe convergence. It sees only the starting
distribution of agent confidences and critiques, which encodes a snapshot of
prior agreement but not the response-to-critique dynamics that generate the
downstream signal. This is a structural limitation of the minimal probe we
test, not of the framework.

\paragraph{Implications for upstream systems.} For
\textsc{CellForge}~\citep{chen2025cellforge} and
\textsc{MassGen}~\citep{massgen2025github}, the immediate actionable
recommendation is: \emph{emit the $\{C(r), S(r), w(r)\}$ trace as a
first-class output} of every orchestrator run. TDI and its diagnostics cost
less than $1$\,ms per run and are trivial to surface through an existing
skill interface; even without a calibrated recommender, the run-level
quantities (ACE trajectory, $\Delta C$, WFR) are directly useful to human
operators as a sanity check.

% =====================================================================
\section{Limitations and Threats to Validity}\label{sec:limits}

We flag five limitations that a reviewer should weigh.

\paragraph{Synthetic DGP.} Every empirical result in this paper comes from a
hand-designed data-generating process (\S\ref{subsec:dgp}). The DGP is
intentionally simple and its parameters are reported in full, but at least
three of its assumptions may fail on real runs:
(i) confidence rates rise linearly in $r$;
(ii) critique severities couple to target confidence with a fixed weight;
(iii) the utility model is a saturating exponential in $N\!\cdot\!R\!\cdot\!b_f$.
Real LLM-driven agent dynamics may be noisier, more bursty (revision spikes),
or influenced by prompt engineering in ways the DGP does not model. We
explicitly did \emph{not} calibrate DGP parameters against measured
\textsc{CellForge} or \textsc{MassGen} runs; doing so is the single most
important follow-up experiment.

\paragraph{Probe weakness.} Our minimal probe is one shallow round. Its
four-dimensional signature predicts difficulty at $R^2\!=\!0.16$ on held-out
tasks --- enough to tilt the prior but not to strictly dominate uniform
allocation at large budgets. Richer probes are natural:
(a) a \emph{two-round} shallow probe that admits a primitive $\Delta C$
estimate; (b) a \emph{per-agent calibration} that down-weights historically
over-confident agents via a Brier-style reliability prior; (c) a
\emph{learnable probe} (MLP or tiny transformer) that maps a full trace
segment to an implicit difficulty embedding. We chose the minimal probe
because it sets a conservative floor.

\paragraph{Configuration space.} $|\Phi|\!=\!12$ here. Larger $\Phi$ would
stretch the Gaussian diagonal-covariance likelihood we use; a Gaussian mixture
or a neural density estimator is a straightforward next step. Our Bayesian
recommender is also a \emph{MAP} estimator; full posterior sampling would
expose uncertainty that is currently swept into point predictions.

\paragraph{Reward model free-riding.} We deliberately do not train a process
reward model. A PRM over agent traces would likely outperform our
hand-designed metrics, but it would also require an expensive labelling phase.
The paper's claim is that hand-designed difficulty metrics over the
already-produced trace are cheap and already actionable; a PRM is a
complementary rather than competing direction.

\paragraph{Generalisation across problem domains.} Our claims specialise to
\emph{convergence-shaped} problems: propose-critique-vote systems where the
team's confidence distribution meaningfully updates round over round. In
problem domains where one agent simply answers ``I know'' with high confidence
regardless of difficulty (e.g.\ rote retrieval), the TDI signal will be
degenerate and the probe uninformative. The framework is not designed to
replace a domain-specific accuracy evaluator.

\paragraph{Reproducibility.} All experiments use \texttt{numpy.default\_rng(2026)}.
The CSV artifacts, figures, and this paper's \LaTeX\ source are produced by
three scripts in \texttt{paper/experiments/} that run in under a minute on a
CPU without any network access.

% =====================================================================
\section{Reproducibility}\label{sec:reproducibility}

Every number, figure, and table in this paper is produced by a three-script
pipeline that runs in under 10 seconds on a single CPU without any network
access. The pipeline is:
\begin{itemize}[leftmargin=*,topsep=0pt,itemsep=0pt]
  \item \texttt{paper/experiments/simulate.py} --- runs seven experiments
        (E1, E1b, E2, E2b, E3, E4, E5) and emits six CSVs into
        \texttt{paper/experiments/out/}.
  \item \texttt{paper/experiments/plot.py} --- renders the five figures as
        PDF\,+\,PNG from the CSVs.
  \item \texttt{paper/experiments/generate\_tables.py} --- emits the five
        \texttt{.tex} snippets the manuscript \texttt{\string\input}s.
\end{itemize}

\paragraph{Seeded determinism.} All randomness flows through
\texttt{numpy.random.default\_rng(2026)}, instantiated once in
\texttt{simulate.py::main}. Rerunning the script with the same seed
reproduces the CSVs byte-for-byte on any platform with identical
\texttt{numpy} ABI behaviour.

\paragraph{DGP parameters are frozen in-source.} All $17$ parameters of the
data-generating process (\S\ref{subsec:dgp}) are declared in a single
\texttt{DGP} dict at the top of \texttt{simulate.py}. Changing one value and
rerunning yields a clearly diffable delta in every downstream artifact.

\paragraph{Artifact dependency graph.}
\begin{center}\small\ttfamily
simulate.py $\to$ out/*.csv $\to$ \{plot.py, generate\_tables.py\}\\
$\to$ figures/*.\{pdf,png\} $+$ tables/*.tex $\to$ paper.tex $\to$ paper.pdf
\end{center}
No file in the dependency graph is generated outside this pipeline.

\paragraph{Software stack.} The entire pipeline requires only
\texttt{numpy}, \texttt{matplotlib}, and a TeX Live distribution with
\texttt{natbib}, \texttt{booktabs}, \texttt{hyperref}, \texttt{caption}, and
\texttt{enumitem}. We do not depend on \texttt{scipy},
\texttt{scikit-learn}, or \texttt{pandas}. The \texttt{algorithm} /
\texttt{algpseudocode} packages are deliberately avoided (some TeX
distributions omit them) --- \S\ref{sec:probe}'s procedure is rendered as a
plain numbered list inside an \texttt{\string\fbox} so compilation succeeds
on a minimal TeX install.

\paragraph{Tests.} Forty-three unit tests under \texttt{tests/} cover the
metric, probe, recommender, instrumentation, calibration, and model modules.
They execute in under 100 ms.

\paragraph{Sixteen tests for the reference TTC implementations.} The
framework-free reference implementations of best-of-$N$, majority vote,
weighted majority, iterative revision, and adaptive-budget allocation in
\texttt{research/test-time-compute-guide/ref\_impl/} are covered by their own
16-test suite; these are pedagogical rather than experimental but we include
them so the single-chain TTC comparisons in \S\ref{sec:related} are not
claims over a reader's prior knowledge alone.

\paragraph{Companion design document.} The file \texttt{docs/DESIGN.md} in
the project records the contribution audit, the MassGen bridge specification,
and a reviewer-oriented validation checklist. It also lists three
falsifying outcomes that would contradict the paper's central claims and
explains exactly which one-line configuration change surfaces each.

% =====================================================================
\section{Conclusion}\label{sec:conclusion}

We have proposed ACE, CSD, $\Delta C$-based convergence signals, and WFR as a
four-feature lens on multi-agent group-generation traces, combined into a
calibrated Task Difficulty Index with Spearman $\rho\!=\!0.92$ against a
ground-truth difficulty label. We have positioned the TDI and its probe as
the multi-agent analogue of \citet{snell2024scaling}'s compute-optimal
single-chain allocation, and released a Bayesian recommender that turns
agentic orchestration into a small and well-posed hyperparameter-tuning
problem. Our experiments are synthetic, we flag the probe-weakness limitation
honestly, and we outline the specific empirical validation on real
perturb-seq data with an \textsc{scGPT} backbone that would move the claim
from ``framework works on a documented DGP'' to ``framework routes production
orchestration compute usefully''.

\paragraph{Code and data.} All artifacts are released under Apache-2.0 at the
repository that accompanies this submission, including the
\textsc{MassGen}-installable skill that produces per-run metric reports for
any propose-critique-vote orchestrator.

% =====================================================================
\small
\bibliographystyle{plainnat}
\bibliography{references}

\end{document}
